\subsection{A fixed-local tripod obstruction inside every restricted product}

\begin{statement}{A.142}{Proved}
{An ordinary completed pair whose new coordinate fails quadratic reciprocity}
Put
\[
 \begin{gathered}
 h=\frac12,\qquad R=\frac{529}{382},\qquad
 c=-\frac{338}{191},\qquad c'=1-c=\frac{529}{191},\\
 b=4c(1-c)=-\frac{715208}{36481},\qquad
 d=18512174429,\qquad p_*=4679,\\
 T_* =c'+p_* =\frac{894218}{191},\qquad
 B_0=(8\cdot128+8)^2=1065024,\\
 S_0=\{p:p\le B_0\}\cup\{1119169\}.
 \end{gathered}
\tag{A.142.1}
\]
There is an ordinary completed unit-equation pair
\((A,B)\in(\widehat{\Q^\times})^2\), with infinitely many
nonzero ordinary integer valuations in each coordinate, and fixed
centres
\[
 x_p\in\Q_p\setminus\{0,1,b\}\quad(p<\infty),\qquad
 A_p=\kappa_p(x_p),\quad B_p=\kappa_p(1-x_p).
\tag{A.142.2}
\]
The real signs are those of the ordinary centre
\(x_\infty=4R(1-R)=-77763/36481\).
The new ordinary values
\[
                    Y_p=x_p-b
\tag{A.142.3}
\]
are units for every \(p\notin S_0\).  Thus their local Kummer
classes belong to the finite-coefficient restricted product at
every exponent.  Nevertheless the actual quadratic evaluations are
\[
 \operatorname{inv}_p(Y_p,d)=
 \begin{cases}
  1/2,&p=p_*,\\
  0,&p\ne p_*,
 \end{cases}
 \qquad \operatorname{inv}_\infty(x_\infty-b,d)=0.
\tag{A.142.4}
\]
Consequently no rational squareclass has these localizations at all
finite primes, and no completed new coordinate has the fixed local
values \(x_p-b\).  Equivalently, there is no rational coefficient
\(e\in\Q^\times\) for which every prescribed point \(x_p\) lifts
to \(t-b=eZ^2\).

This constructs fixed local data satisfying the restricted-product
condition of A.132 and failing its character-orthogonality condition.
All local centres and the completed pair are constructed together;
the assertion does not replace the fixed centres by new points when
testing the additional coordinate.
\end{statement}
\begin{proof}
\emph{The exact initial cover and the one changed place.}
Use the actual quartic cover of A.134,
\[
 \begin{gathered}
 C:\quad T-i=a_iX_i^4\quad(i=0,1,h),\\
 a_0=\frac{529}{382},\quad a_1=\frac{147}{382},\quad
 a_h=\frac{4089852535355054689}{191}.
 \end{gathered}
\tag{A.142.5}
\]
Keep its retained ordinary points at all primes at most \(4500\)
and at \(1119169\), and keep its real point over \(R\) with root
signs \((-1,+1,+1)\).  The retained independent certificate verifies
that none of these \(611\) centres is \(c'\), and that
\[
 \operatorname{inv}_p
       \bigl(4T_p(1-T_p)-b,d\bigr)=0
                   \quad\text{at all these places}.
\tag{A.142.6}
\]
One can check this table from the following three literal symbols:
\[
 \begin{aligned}
 4T(1-T)-b&=-4(T-c)(T-c'),\\
 (4T(1-T)-b,d)&=(-4,d)+(T-c,d)+(T-c',d).
 \end{aligned}
\tag{A.142.7}
\]
On the retained centres the second term is nonzero exactly at two;
the third is nonzero exactly at \(2,7,139\); and the constant first
term is nonzero exactly at \(7,139\).  The three lists cancel
place by place.  The certificate checks the exact valuations and
residue units, rather than just this parity count.

At the additional prime \(p_*\), choose \(T=T_*\).
The integer \(4679\) is prime and is \(23\) modulo \(24\).
At the boundary \(c'\) the three literal radicands are
\[
 \frac{c'}{a_0}=2,\qquad
 \frac{c'-1}{a_1}=\frac{676}{147},\qquad
 \frac{c'-h}{a_h}
        =\frac{867}{2(13\cdot155564491)^2}.
\tag{A.142.8}
\]
They have respective fourth roots
\(2304,2961,4466\) in \(\F_{4679}^\times\), as checked by
literal exponentiation.  Since \(T_*\equiv c'\pmod{p_*}\), these
are also residue roots of the three equations at \(T_*\).
Their derivatives \(4X_i^3\) are units, so each root has a unique
actual Hensel lift to the indicated residue branch in \(\Q_{p_*}\).
Thus this is an ordinary point of the actual cover, with all three
root coordinates units.

The number \(d\) is a nonsquare unit at \(p_*\), and
\[
       T_*-c'=p_*,\qquad T_*-c\in\Z_{p_*}^\times,
          \qquad v_{p_*}(4T_*(1-T_*)-b)=1.
\tag{A.142.9}
\]
The associated quadratic extension is unramified and its norm
group has even valuations.  This proves the value \(1/2\) at
\(p_*\) in (A.142.4).

\emph{All remaining initial controlled points.}
At every other prime \(4500<p\le B_0\), take an ordinary
unit-coordinate residue point of the old curve \(C\).
It exists by the genus-33, old-boundary-64 estimate
(A.134.13); every such prime is a good prime of (A.142.5), apart
from the separately retained prime \(1119169>B_0\).
If its centre reduces to \(c\), use the rational lift
\(T=c+p^2\); if it reduces to \(c'\), use \(T=c'+p^2\).
Otherwise use a rational lift in its residue class.
The reductions of \(c,c'\) are distinct here, since
\[
                         c'-c=867/191
\tag{A.142.10}
\]
has support only at \(3,17,191\).
All coordinate roots remain units and lift by their unit
derivatives.  Each lift avoids the old punctures, \(c\), and
\(c'\).  At these primes the valuation of the right side of
(A.142.7) is either zero or two, while \(d\) is an unramified
unit character.  Its invariant is therefore zero.

This prescribes ordinary points at every prime of the finite set
\(S_0\), including the altered point at \(p_*\).
It retains precisely the two dyadic/real pairs of nonzero initial
root symbols in (A.134.21).  Indeed all additional roots are units,
and the fixed boundary fields
\[
 \begin{gathered}
 B_0^{\rm bdry}=\Q(\sqrt{-1146},\sqrt{-382}),\qquad
 B_1^{\rm bdry}=\Q(\sqrt{382}),\\
 B_h^{\rm bdry}=\Q(\sqrt{191},\sqrt{-573})
 \end{gathered}
\tag{A.142.11}
\]
are unramified at all these additional primes.
Their unit symbols are zero.  In particular the changed point
at \(p_*\) introduces no root-symbol obstruction to the old
coordinate refinement.

\emph{The modified mixed-cover induction.}
We now retain the displayed initial points literally and apply
the exact refinement construction of A.134, with the following
extra open condition: both \(T-c\) and \(T-c'\) are units
outside the controlled set, and both are units at every newly
introduced coefficient prime.  We verify this additional condition
and all the inputs that permit the construction.

For mixed exponents \(\boldsymbol n=(n_0,n_1,n_h)\), all divisible
by four, put \(D=n_0n_1n_h\) and
\(N=\operatorname{lcm}(n_0,n_1,n_h)\).
The geometric degree, inertia and genus are the actual ones of
(A.90.18).  Removing both unramified fibres over \(c,c'\) gives
\[
 \begin{aligned}
 2g-2&=-2D+\sum_i D(1-1/n_i)+D(1-1/N),\\
 b_{\rm geom}&=\sum_i D/n_i+D/N+2D\le6D.
 \end{aligned}
\tag{A.142.12}
\]
Thus the same explicit bound
\[
                         B(\boldsymbol n)=(8D+8)^2
\tag{A.142.13}
\]
gives \(p+1-2g\sqrt p-b_{\rm geom}>0\) for every
\(p>B(\boldsymbol n)\).  Indeed \(2g\le2D+2\) and
\(\sqrt p>8(D+1)\) make this inequality strict even with
the bound \(6D\).  At good primes the resulting residue point
has all roots nonzero and both displayed differences units;
unit-derivative Hensel lifting preserves these properties.
The first refined exponent vector is \((8,4,4)\), for which
(A.142.13) is the number \(B_0\) in (A.142.1).
At the outset, every prime outside \(S_0\) exceeds this bound
and also the smaller bound for \((4,4,4)\); the same calculation
therefore supplies the initial outside-set points with both
differences units.

All initial coefficient, exponent, boundary-difference and
boundary-field ramification primes lie in \(S_0\).
This includes the differences \(i-c'\): their supports are
contained in \(\{2,3,13,17,23,191\}\).
No field in (A.142.11) contains \(i=\sqrt{-1}\).
Coefficient changes by the current even exponent preserve their
squareclasses.  Consequently the actual mixed boundary common
fields remain exactly (A.142.11), by (A.90.3)--(A.90.10).

The first refinement is the signed one of (A.134.23)--(A.134.25).
Its initial root-symbol sums are unchanged, as checked above.
For its coordinate-zero target \(x_p=X_0(P_p)\), the negative
rational number
\[
             d_0=-\prod_{p\in S_0}p^{v_p(x_p)}
\tag{A.142.14}
\]
subtracts the principal idele and makes the exact CRT automorphism
fix \(B_0^{\rm bdry}\).  The real quotient \(x_\infty/d_0\)
is positive.  The actual boundary root stabilizers therefore
provide the prime product giving a negative multiplier \(u\).
The literal map is
\[
             X_0=uY_0^2,\qquad a_0'=a_0u^4,\qquad n_0'=8,
                    \qquad T=T,\quad X_1=Y_1,\quad X_h=Y_h.
\tag{A.142.15}
\]
At the controlled places, sufficiently close approximation gives
a principal-square root \(Y_0\); at the real place choose its
positive root.  This establishes the individual zero-symbol
invariant, exactly as in A.134.

Every later refinement of coordinate \(i\) by an integer \(m\)
uses the actual maps
\[
   X_i=uY_i^m,\qquad a_i'=a_i u^{n_i},\qquad n_i'=mn_i,
               \qquad T=T,\quad X_j=Y_j\ (j\ne i).
\tag{A.142.16}
\]
The individual zero symbols give the precise ray-class condition
in (A.90.20)--(A.90.23).  Its proved stabilizer-image equality
and Chebotarev construct the multiplier, with one extra completely
split prime of valuation one.  The choice is positive at the real
place after (A.142.15).  All controlled centres are retained
exactly, and the new roots are principal local squares.

At a fresh coefficient prime \(q\), the construction takes
\(Y_i=1\) and \(T=i+a_i u^{n_i}\).
The other roots lift the fixed boundary residue tuple.
The prime is chosen away from the supports of \(i-c\) and
\(i-c'\), whereas \(v_q(T-i)>0\).  Hence both \(T-c\)
and \(T-c'\) are units there.  All new root coordinates are
units and every boundary-field character is unramified, so their
symbols vanish.  At every other prime outside the new controlled
set, (A.142.12)--(A.142.13) give a unit-coordinate point with
both differences units.  When a new prime is subsequently
controlled, retain this point.  These are all the local cases in
the induction.  They prove the stronger permanent assertion
\[
                   T_p-c, T_p-c'\in\Z_p^\times
                                      \quad(p\notin S_0).
\tag{A.142.17}
\]

Refine the other two coordinates by two after the first signed
step, then use cycles with factors \(3,4,5,\ldots\).
The common exponent after the factor \(m\) cycle is \(4m!\).
The coefficient identities in (A.142.16) therefore give compatible
classes \(A_0,A_1,A_h\) in the full multiplicative completion,
with \(A_i\bmod4=[a_i]\).
Every prime is eventually controlled and its centre is then
retained literally.  The late equations with exponent divisible
by an arbitrary prescribed modulus prove
\[
                   (A_i)_p=\kappa_p(T_p-i).
\tag{A.142.18}
\]
Each coordinate has infinitely many persistent fresh integer
valuations, by its extra completely split prime at every
refinement.  Subsequent multipliers are units at that already
controlled prime.  This is the same exact support argument as
in A.134, and none of the extra open conditions changes it.

\emph{The actual tripod pair and the failed global test.}
Define
\[
       A=-4A_0A_1,\qquad B=4A_h^2,
                       \qquad x_p=4T_p(1-T_p).
\tag{A.142.19}
\]
The polynomial identities
\(1-4T(1-T)=4(T-h)^2\) and (A.142.7) prove all the
localization identities and show that \(x_p\) avoids
\(0,1,b\).  This is the actual degree-two map of A.107;
on each finite old cover its root-coordinate map is
\(X=X_0X_1\), \(Y=X_h^2\), with coefficients
\(-4a_0a_1\) and \(4a_h^2\).
The disjoint extra primes for coordinates zero and one give
infinite support of \(A\); those for coordinate \(h\) give
twice their nonzero valuations in \(B\).
The real centre remains \(R\), giving the signs asserted in
the statement and
\(x_\infty-b=637445/36481>0\).

All primes of \(S_0\) retain their initial centres.
Equations (A.142.6), (A.142.9), and the prescribed even-valuation
lifts at the other initial good primes prove (A.142.4) on that
finite set.  At every prime outside \(S_0\), (A.142.17) makes
\(Y_p=-4(T_p-c)(T_p-c')\) a unit, and \(d\) is an unramified
unit character.  Its invariant is zero.  At the real place
\(d>0\), so the character is split.  This proves the full
invariant statement at all places.

For every integer \(N\), the tuple
\((\kappa_N(Y_p))_p\) lies in the restricted product of
\(H^1(\Q_p,\mu_N)\) with respect to the unramified subgroups
away from \(N\): all places outside \(S_0\cup\{p:p\mid N\}\)
are units at primes not dividing \(N\).
For \(N=2\), however, its cup pairing with the actual global
character of \(\Q(\sqrt d)/\Q\) has sum \(1/2\).
If a rational squareclass \([e]\) had the asserted finite
localizations, the global class \((e,d)\in\operatorname{Br}(\Q)[2]\)
would have these local invariants.  Its real invariant is zero
because \(d>0\).  Global reciprocity
\cite[Theorem 12.1.8]{AgentCTSkoro2019} would make their sum zero,
a contradiction.  Projection modulo two excludes every completed
new coordinate as well.
The equation \(t-b=eZ^2\) has a point above a fixed \(x_p\)
exactly when \([e]=[Y_p]\), proving the stated geometric
equivalence with its actual cover map.

\emph{Choice convention and execution scope.}
Keep all \(611\) retained centres and the three special dyadic
root signs of A.134.  Use the specified point and residue roots
at \(p_*\).  At other initial good primes choose the first
ordinary residue point in a fixed lexicographic order and apply
the explicit centre-lifting rule above.  At later stages choose
the first residue point in the open curve with both new fibres
removed.  Use the least sufficient Hensel/conductor precisions,
the fixed enumeration of actual mixed root-stabilizer words,
and the least admissible distinct Frobenius primes and extra
split prime, as in A.134's effective convention.
Each finite selection terminates by the point-count and
Chebotarev arguments already proved.  These conventions specify
one completed tuple.  The retained certificate executes the
initial finite arithmetic, not the infinite recursion or the
full initial list of primes up to \(B_0\).
\end{proof}

The exact certificate
\path{code/tripod_reciprocity_obstruction_certificate.py}
and its retained output check the two curve identities, all
\(611\) original local evaluations, the prime \(4679\), its
literal fourth roots, the new nonzero symbol, the complete
support exclusions, the initial root-symbol compatibility,
and the mixed-cover bound with both removed fibres.
